% Appendix : the weak-coupling decomposition theorem on which the localised-stabilisation design rests.
\section{The weak-coupling decomposition theorem (prior work)}
\label{app:theorem}

This appendix states the weak-coupling decomposition theorem on which the localised-stabilisation
design of Section~\ref{sec:localstab-methods} rests. The theorem and its proof are the contribution
of the earlier BBCN118 preprint~\cite{bhatti2026} (CC BY 4.0); the statement and proof sketch below
are reproduced and adapted from that work, so that Section~\ref{sec:localstab-methods} can be read
without consulting the source. The remark at the end is the new content. It records that the theorem
fails on the pathway partition but holds on the timescale partition, which is the move the paper
makes.

\paragraph{Setup.} Let a Boolean network with global update $x(t{+}1)=F(x(t))$ be partitioned into
$m$ pathway modules $P_1,\dots,P_m$, with module states $x_i$ and module targets $x_i^{*}$, so that
$x=(x_1,\dots,x_m)$ and $x^{*}=(x_1^{*},\dots,x_m^{*})$.

\paragraph{Theorem 1 (convergence under pathway decomposition).} Assume:
\begin{itemize}
\item[(A1)] \emph{Weak inter-pathway coupling:} $\displaystyle\max_{i\neq j}\left\|\partial F_i/\partial x_j\right\| < \varepsilon$, with $\varepsilon \ll 1$;
\item[(A2)] \emph{Local kernel effectiveness:} for each module $i$ there is a kernel $u_i$ that drives $x_i(t)\to x_i^{*}$ in finitely many steps; and
\item[(A3)] \emph{Modularity:} $F_i(x)\approx f_i\!\left(x_i, e_i(x)\right)$, where $e_i(x)$ collects the inputs $P_i$ receives from other modules.
\end{itemize}
Then the sequential application of the pathway kernels $u=\bigcup_{i=1}^{m} u_i$ drives the global
state into a neighbourhood of the target,
\[
x(t)\to \mathcal{N}_\delta(x^{*})=\{x : d_H(x,x^{*})\le \delta\},\qquad \delta = O(\varepsilon),
\]
where $d_H$ is the Hamming distance.

\paragraph{Proof sketch.} Take the weighted Hamming Lyapunov function
\[
V(x)=\sum_{i=1}^{m} w_i\, d_H\!\left(x_i,x_i^{*}\right),\qquad w_i>0 .
\]
$V$ is non-negative, integer-valued, bounded above by $\sum_i w_i|P_i|$, and zero exactly at $x^{*}$.
Intervening on module $k$ reduces its term by some $\alpha_k>0$ by (A2). By (A1) every other module
moves by at most $w_j\varepsilon|P_j|$, so
\[
\Delta V \le -\alpha_k + \varepsilon\sum_{j\neq k} w_j|P_j| ,
\]
which is strictly negative once $\varepsilon$ is small enough relative to $\alpha_k$. Because $V$ is
bounded below, strictly decreasing, and integer-valued, it reaches a minimum $V_{\min}$ in finitely
many steps, with $V_{\min}\le\delta=O(\varepsilon)$; at $\varepsilon=0$ the bound gives $V_{\min}=0$
and exact convergence. \hfill$\square$

\paragraph{Corollary 1 (convergence rate).} If each pathway kernel reduces the mismatch by at least
$\beta>0$ per intervention, the system reaches the neighbourhood in at most $T \le V(x(0))\,/\,(\beta/2)$
interventions.

\paragraph{Remark: the partition is what matters.} The result holds only under assumption (A1). On
the pathway partition the present network violates it. Its pathway-coupling graph carries a dominant
strongly-connected component of 16 of the 22 regulatory pathways (Section~\ref{sec:obstruction}), so
the inter-module coupling is not small and pathway-wise stabilisation does not deliver joint
reachability. The paper does not abandon the theorem. It re-partitions the apoptosis core by process
timescale, into fast, mid, and slow blocks. The delay embedding makes each block quasi-static over the
others' update intervals, so the inter-block coupling satisfies (A1) where the inter-pathway coupling
did not, and the theorem applies block by block (Section~\ref{sec:localstab-methods}). The full proof,
including the explicit $\varepsilon$ threshold, is in the cited preprint and is not reproduced here.